% HD-N08--N09. Author: GPT-6 Astra Pro, /root/imprimitive.
% 21 November 2026. Reasoning effort inherited; numerical setting not exposed.
% Geometry and quantitative normalization developed jointly with /root.
\section{An effective coherent CM source for every class-number-one field}

\paragraph{HD-N08 (two bielliptic directions, an actual sheaf, and regularity).}
Let $d$ be a positive squarefree integer and put
\[
 K=\mathbb Q(\sqrt{-d}),\qquad D=4d,\qquad k=D+1,\qquad r=64D.
\]
The construction below gives a perfect coherent sheaf $F$ of rank $r$
on the actual CM abelian sixfold $A=E_i^3\times\widehat{E_i^3}$,
with projective dimension at most one. It is strongly semiregular for
an explicitly transported finite translation/tensor action. Its
balanced $K$-structure is the one in HD-N07 with parameter $D$.
For
\[
 \mathcal H=3(H_x+H_\alpha)
\]
it is $560$-regular. More precisely, for every integer $m\ge1$ and
every line bundle $M$ such that $c_1(M)-560m^2\mathcal H$ is nef,
\[
 H^j(A,[m]^*F\otimes M)=0\qquad(j>0).                 \tag{N8.1}
\]
The field need not have class number one for this source construction.

\paragraph{The two actual quotient ideals.}
Use the Fermat plane quartic $C$, its three bielliptic maps
$f_i:C\to E_i$, and the degree-eight isogeny
\[
 q:X=E_i^3\longrightarrow J(C),\qquad q^*\Theta=2H,
 \qquad H=\eta_1+\eta_2+\eta_3
\]
proved in HD-N01. Its kernel lies in $X[2]$. Write
$U_i=f_i^*E_i\subset J(C)$. The exact local-inertia proof of
(N7.1) gives
\[
 U_i\cap(C-C)=\{0\}\qquad(i=1,2).                    \tag{N8.2}
\]
Choose points $\tau_1,\tau_2$ of exact order $k$ on the first and
second elliptic factors, respectively. For $i=1,2$, let
$\epsilon_i=f_i^*\tau_i$, let
$u_i:J(C)\to Y_i=J(C)/\langle\epsilon_i\rangle$, and let
$\overline C_i=u_i(C)$. Equation (N8.2) says that every orbit
consists of $k$ pairwise disjoint translates, so $u_i|_C$ is a
closed embedding. Put
\[
 r_i=u_iq,\quad K_i=\ker r_i,\quad
 F_i=r_i^*I_{\overline C_i}=I_{Z_i}.
\]
Thus $|K_i|=8k$, and the two kernels are explicitly
\[
 K_1=\ker q+\langle(\tau_1,0,0)\rangle,\qquad
 K_2=\ker q+\langle(0,\tau_2,0)\rangle.                \tag{N8.3}
\]
Their quotient ideals have exactly the Ext and trace calculation
of HD-N02 and HD-N07: $\operatorname{Ext}^1$ and
$\operatorname{Ext}^2$ each have dimension six; the degree-two
trace is injective, as are the degree-one and degree-zero traces.
The calculation uses the isomorphism
$H^1(\mathcal O_{Y_i})\to H^1(\mathcal O_C)$ and the canonical
multiplication isomorphism
$\operatorname{Sym}^2H^0(K_C)\to H^0(K_C^2)$ for a plane quartic.
The coefficient of the quadratic polarization term is $-k/2$,
which is nonzero. In particular the same two trace blocks have
ranks three and three.

The external product of the two quotient ideals is semiregular:
the three Kunneth summands of its degree-two self-Ext group map
injectively to the separate Hochschild bidegrees $(2,0)$, $(1,1)$
and $(0,2)$. Pullback by $r_2\times r_1$, reflection in the first
factor, and tensoring by $\mathcal O_X(2H)$ in each factor give
the corresponding strong equivariant semiregularity. This uses
two different quotient ideals; the trace proof does not require
them to be the same sheaf.

\paragraph{The Fourier transform and the section-injectivity proof.}
Let $L=\mathcal O_X(2H)$, $\mu(x,y)=(x+y,y)$, and use the inverse
Fourier kernel $\mathcal P^{-1}[3]$ in the second variable. Define
\[
 E_{\mathrm{ref}}=
 (\operatorname{id}\times\mathrm{FM}_{\mathcal P^{-1}[3]})
 \mu^*\bigl(([-1]^*F_2\otimes L)\boxtimes(F_1\otimes L)\bigr).
                                                               \tag{N8.4}
\]
The order of the two factors makes the following incidence
description literal. On $X_y\times A$, put
\[
 a(y,x)=y+x,\quad
 \mathcal L=L_y\otimes a^*L\otimes\mathcal P_{y,\alpha}^{-1},
 \quad D_1=\{y\in Z_1\},\quad D_2=\{y+x\in-Z_2\}.
\]
For $\pi:X_y\times A\to A$ and $T=D_1\cap D_2$, set
\[
 G=R\pi_*(\mathcal L I_{D_1}I_{D_2}),
 \qquad E_{\mathrm{ref}}=G[3],\qquad F=G^\vee[-1]
               =E_{\mathrm{ref}}^\vee[2].                 \tag{N8.5}
\]
Every equality here is an actual derived-functor identification.

Here is the additional geometric fact that makes $F$ a sheaf.
The curve $C'=q^{-1}C$ is connected: the Abel map induces the
surjection $\pi_1(C)\to\pi_1(J(C))$, hence has transitive monodromy
on the fibre of the connected etale cover $q$. It is an etale
degree-eight cover of $C$, and therefore has genus17. Projection
of $C'$ along the first elliptic factor is injective. Indeed, two
points in one such fibre differ by a first-factor point; their
images in $C$ differ by an element of $U_1$, so (N8.2) makes those
images equal. The difference then lies in
$\ker q\cap(E_i\times0\times0)=0$, since $f_1^*$ is injective.
The same argument applies to the second factor.

The two projected images are irreducible divisors $B_1,B_2$ on
the corresponding elliptic surfaces. Their projection maps have
degree one, so the actual inverse-image divisors $S_1,S_2$ on $X$
have classes
\[
 [S_1]=4(\eta_2+\eta_3),\qquad
 [S_2]=4(\eta_1+\eta_3).                                  \tag{N8.6}
\]
This follows by pushing forward
$[C']=q^*(\Theta^2/2)=2H^2$ along each elliptic projection.
There is no unaccounted projection degree. One may also check
smoothness: the projected divisor has self-intersection32 and
arithmetic genus17, equal to the genus of its degree-one
normalization $C'$, so the normalization has no singularity defect.

The $k$ translates defining $Z_1$ lie on $S_1$. On every generic
first-factor fibre of $S_1$ they supply $k$ distinct points.
The $k$ reflected translates in $-Z_2-x$ have the analogous
property on the second-factor fibres of $-S_2-x$. A section of
$\mathcal O_X(4H)\otimes P$ vanishing on both unions, where
$P\in\operatorname{Pic}^0(X)$ is arbitrary, must vanish on each
of these entire elliptic fibres: its restriction has degree four,
whereas $k=4d+1\ge5$. Hence its divisor contains both $S_1$ and
$-S_2-x$. They are distinct for every $x$, because their classes
in (N8.6) differ. The residual effective divisor would have class
\[
 4H-[S_1]-[S_2]=-4\eta_3.
\]
Its intersection with the ample class $H^2$ would be $-8$, which
is impossible for an effective divisor. Thus that space of
sections is zero, for every $x$ and every flat factor $P$.

\paragraph{An explicit derived-to-coherent calculation.}
The coordinates $(y,z=y+x)$ identify the two curve-union ideals
with pullbacks from independent factors. They are Tor independent,
and their product equals their intersection. Therefore
\[
 0\to\mathcal L I_{D_1}I_{D_2}\to\mathcal L
 \to\mathcal L|_{D_1}\oplus\mathcal L|_{D_2}
 \to\mathcal L|_T\to0                                    \tag{N8.7}
\]
is exact on the total space.
On a fibre of $\pi$, $\mathcal L$ has class $4H$, with
$h^0=64$ and no higher cohomology. Its restriction to a component
$C'$ has degree48. Since $g(C')=17$, its first cohomology is zero
and its space of sections has dimension32. Consequently
\[
 W=\pi_*\mathcal L,\qquad
 U=\pi_*(\mathcal L|_{D_1}\oplus\mathcal L|_{D_2})
\]
are vector bundles of ranks64 and $64k$, respectively.

The map $f:T\to A$ is finite flat over each of its smooth divisor
images. To see the precise degree, ignore translations and write
one component as $C'\times C'\times\widehat X$, with its map to
$X$ given, up to sign, by addition. The diagonal action
$(p',q')\mapsto(p'+h,q'-h)$ of $\ker q$ is free. Its quotient is
the base change of $C\times C\to C^{(2)}\simeq\Theta$ by $q$.
The quotient map is etale of degree eight; the latter map is
finite flat of degree two onto the smooth divisor $q^{-1}\Theta$.
Thus the composite has degree16. Smoothness of $\Theta$ follows
from the absence of a degree-two pencil on the nonhyperelliptic
genus-three curve. Each summand of
\[
 \mathcal T=f_*(\mathcal L|_T)
\]
is a vector bundle on a smooth divisor, pushed forward to $A$.
It has projective dimension one.

Applying direct image to (N8.7), including the preceding higher
cohomology vanishings, represents $G$ by
\[
 [W\longrightarrow U\longrightarrow\mathcal T]
 \quad\hbox{in degrees }0,1,2.                            \tag{N8.8}
\]
The section argument proves that $W\to U$ is fibrewise injective.
Hence $Q=U/W$ is a vector bundle of rank $64(k-1)=r$, and
$G=[Q\to\mathcal T]$ in degrees1,2. No assertion that fibrewise
$H^2(G)$ vanishes has been made or used.

Finite duality gives
\[
 \mathcal T^D=\mathcal E xt^1_A(\mathcal T,\mathcal O_A)
 =\bigoplus f_{ab,*}
       ((\mathcal L|_{T_{ab}})^{-1}\otimes\omega_{T_{ab}}),
\]
with all other sheaf Ext groups zero. Dualizing the last two-term
complex and shifting by $-1$ gives an ordinary coherent sheaf and
the exact extension
\[
 0\longrightarrow Q^*\longrightarrow F
   \longrightarrow\mathcal T^D\longrightarrow0.          \tag{N8.9}
\]
Both the local extension and its finite-duality description prove
$\operatorname{pd}F\le1$; local freeness outside a prescribed
codimension is unnecessary.

\paragraph{The actual action, trace and whole character.}
In (N8.4), the first translation parameter belongs to $K_2$ and
the second to $K_1$. With the pullback convention of HD-N07,
the transported action on $E_{\mathrm{ref}}$ has point map
\[
 (x,y)\longmapsto
 \bigl((x-y,2(x+y)),\,(-2x,y)\bigr)
 \quad\hbox{in }A\times\widehat A.                       \tag{N8.10}
\]
It is injective, with group order $64k^2$, and inherits the full
coherence from finite descent. Duality replaces a translation/tensor
point $(a,b)$ by $(a,-b)$; invert the dual linearizations to obtain
the coherent action on $F$. The trace is preserved up to invertible
signs under this operation. Explicitly, derived transpose identifies
the degree-two self-Ext spaces, the Atiyah class of the dual is
the negative transpose, and graded cyclicity of the trace gives
the sign $(-1)^p$ on
$\operatorname{tr}(\xi\operatorname{At}^p/p!)$.
The reversal signs for the cohomological and differential-form
factors cancel. Thus the complete trace kernels identify, including
the invariant subspaces. The even shift $[2]$ contributes sign+1.
This proves the claimed strong equivariant semiregularity of $F$.
Negative self-Ext groups vanish because $F$ is a sheaf; they also
vanish through the displayed equivalence with the original ideals.

Each factor in (N8.4) has the same Chern character as the factor
used in HD-N07. Reflection acts trivially in even cohomology.
Consequently the exact Fourier calculation there, with parameter
$D$, applies without alteration. Put
\[
 \theta_D=4D H_x+H_\alpha,\quad B=-P/2,
 \quad\lambda_D=\frac{D+1}{2D}.
\]
For the primitive integral Weil class $t_D$ of HD-N07 one has
\[
 e^{-B}\operatorname{ch}(F)=-\lambda_Dt_D+P_D(-\theta_D),
                                                               \tag{N8.11}
\]
where the entire polynomial is
\[
 P_D(z)=64D-8(D-1)z-2z^2+\frac{D-1}{12D}z^3
       +\frac{z^4}{96D}-\frac{D-1}{3840D^2}z^5
       -\frac{z^6}{46080D^2}.                              \tag{N8.12}
\]
The rational action of $\sqrt{-D}$ is
$\left(\begin{smallmatrix}0&1/2\\-2D&0\end{smallmatrix}\right)$
on each elliptic pair. Its signature is $(3,3)$, its Rosati
involution is conjugation, and the Hermitian space is the explicit
$U^3$ of HD-N07. The polarization product of $\theta_D$ is $64D^3$.
The class satisfies $\theta_Dt_D=0$ and
$\int t_D^2=-2048D^3$.

\paragraph{The explicit regularity calculation.}
All the following vanishings allow arbitrary flat factors.
Relative duality and Serre duality show, for every integer $q\ge1$,
\[
 H^j(A,W^*(q\mathcal H))=H^j(A,U^*(q\mathcal H))=0
 \qquad(j>0).                                             \tag{N8.13}
\]
Here is the numerical positivity test, so that (N8.13) does not
depend on a qualitative regularity assertion. Project the duality
calculation onto the fixed $y$ coordinate, or onto the fixed curve
coordinate for a summand of $U$. The line bundle on the remaining
six-dimensional fibre has matrix
\[
 \begin{pmatrix}2-3q&0\\0&-3q\end{pmatrix}
 \quad\hbox{or}\quad
 \begin{pmatrix}2-3q&1\\1&-3q\end{pmatrix}.
\]
Both are negative definite: the largest eigenvalue before
subtracting $3qI$ is at most $1+\sqrt2<3$. Only the top direct
image of this antiample line bundle can be nonzero, which gives
the required cohomology vanishing below degree six on the dual
side. This proves (N8.13). In particular $W^*$ and $U^*$ are
$7$-regular.

Take the evaluation sequence
\[
 0\to K_U\to\mathcal O_A(-7\mathcal H)^N\to U^*\to0.
\]
Its kernel is locally free and $14$-regular. For degrees at least
two this follows directly from (N8.13) and positivity of
$(7-j)\mathcal H$. For degree one, the multiplication map from
the degree-seven sections, followed by six hyperplanes in the
embedding defined by $\mathcal H$, is surjective by regularity.
Compose the evaluation map with $U^*\to W^*$ and call the kernel
$J$. Then
\[
 0\to K_U\to J\to Q^*\to0.                              \tag{N8.14}
\]
The Buchsbaum--Rim resolution of $J$ has terms
\[
 C_j=\bigwedge^{64+j}\mathcal O_A(-7\mathcal H)^N
          \otimes\operatorname{Sym}^{j-1}W\otimes\det W,
 \qquad j\ge1.                                           \tag{N8.15}
\]
The first term surjects onto $J$. Local exactness follows by
splitting the surjection as $K\oplus V\to V$ and decomposing
according to the number of factors from $K$. The resulting
strands are the polynomial de Rham complex in positive degree,
contracted by the Euler vector field in characteristic zero;
the degree-one strand maps isomorphically to $K$.

For an explicit bound on (N8.15), the theorem of the square gives
\[
 W=\mathcal O_A(2H_x)\otimes b^*\mathrm{FM}(\mathcal O_X(4H)),
 \quad b=2x-\alpha,
 \quad [4]^*W=Q_0^{\oplus64},
\]
with
\[
 c_1(Q_0)=16H_x-4H_\alpha+8P,\qquad
 \operatorname{Mat}(Q_0)=\begin{pmatrix}16&8\\8&-4\end{pmatrix}.
                                                               \tag{N8.16}
\]
Indeed the pullback of the Fourier bundle by the polarization
isogeny of $\mathcal O(4H)$ is
$H^0(\mathcal O(4H))\otimes\mathcal O(-4H)$; the displayed
formula is its pullback by $b$ with the first line-bundle factor.
The eigenvalues of the matrix in (N8.16) are $6\pm2\sqrt{41}$,
whose minimum is greater than $-7$, while
$[4]^*\mathcal H$ has matrix $48I$. The symmetric and determinant
factor in (N8.15) pulls back to copies of $Q_0^{j+63}$.
Kodaira vanishing and the trace splitting of the isogeny therefore
give
\[
 \operatorname{reg}_{\mathcal H}
   (\operatorname{Sym}^{j-1}W\otimes\det W)\le j+69,
 \qquad \operatorname{reg}C_j\le517+8j.                  \tag{N8.17}
\]
Only six terms are needed to bound $J$ on a sixfold: successive
cohomology sequences leave a final kernel in cohomological
degree seven. Thus
\[
 \operatorname{reg}J\le
 \max_{1\le j\le6}(\operatorname{reg}C_j-j+1)
 \le560,\qquad\operatorname{reg}Q^*\le560.               \tag{N8.18}
\]

For the torsion dual, use $\rho:C'\times C'\to C\times C$.
On $C\times C$ write $F_1,F_2$ for the point fibres,
$S=F_1+F_2$, and $\Delta_C$ for the diagonal. Addition pulls
$\Theta$ back to $A_C=4S-\Delta_C$. The divisor
$R_C=3S-\Delta_C$ is nef. A degree-three map $C\to\mathbb P^1$
from plane projection supplies its effective residual fibre
correspondence. If its monodromy is $S_3$, that residual divisor
is irreducible and has square2, hence is nef. If its monodromy
is cyclic of order three, it is the sum of the two automorphism
graphs: each has square $-4$, and they intersect transversely
in the five fixed points given by Riemann--Hurwitz. Their sum
therefore has intersection1 with each component and nonnegative
intersection with every other curve, again proving nefness.

Put $A'=\operatorname{add}_X^*H=\rho^*A_C/2$ and
$H'_i=\rho^*(3F_i/2)$. Up to flat factors on each translated
component,
$c_1(\mathcal L_T)=2H'_1+2H'_2-P_{1,\alpha}$ and
$f^*\mathcal H=3A'+3H_\alpha$. For $q\ge4$ decompose
\[
 \begin{split}
 qf^*\mathcal H-c_1(\mathcal L_T)
 ={}&[3qA'-3H'_1-2H'_2]+(3q-1)H_\alpha\\
    &+[H'_1+H_\alpha+P_{1,\alpha}].
 \end{split}                                               \tag{N8.19}
\]
The last bracket is a nef addition pullback. The first equals
\[
 \rho^*\left[\frac{3q}{2}R_C+
       \left(\frac{3q}{2}-\frac92\right)F_1+
       \left(\frac{3q}{2}-3\right)F_2\right]
\]
and is ample. Consequently (N8.19) is ample on the smooth
fivefold $C'\times C'\times\widehat X$. Kodaira vanishing and
finite duality give $9$-regularity of $\mathcal T^D$.
Together with (N8.9) and (N8.18), this proves the stated
$560$-regularity of $F$.

For the stronger assertion (N8.1), pull the exact constructions
back by $[m]$. The matrices proving (N8.13) are multiplied by
$m^2$, so the assumed nef difference makes the duality
vanishings hold after tensoring by $M$. The pulled evaluation
sequence has ample middle terms and gives vanishing of
$H^j([m]^*K_U\otimes M)$ for $j\ge2$.
After one further pullback by $[4]$, the first six terms of
(N8.15) are sums of line bundles whose matrices, after removal
of a nef class, are
\[
 m^2\{48(112-7j)I+(j+63)\operatorname{Mat}(Q_0)\},
 \qquad1\le j\le6.
\]
Their least eigenvalue is at least
$m^2(3360-7\cdot69)=2877m^2>0$. The same six-term truncation
and (N8.14) give the desired vanishing for $[m]^*Q^*\otimes M$.
For the torsion term, finite flat base change gives a smooth
etale cover $T'$ of $T$; the required line bundle is the sum
of the pullback of (N8.19) with $q=4$ and a nef class.
Kodaira vanishing on $T'$ proves the last vanishing.
The extension (N8.9) proves (N8.1). \hfill$\square$

\paragraph{HD-N09 (effective bounds on fixed split lattices).}
Let $d$ be one of the integers
\[
 d\in\{1,2,3,7,11,19,43,67,163\},
\]
let
$K=\mathbb Q(\sqrt{-d})$, and let $(A_T,\theta_T)$ be any
polarized split Weil sixfold of signature $(3,3)$, with integral
action of the order of conductor $f$. Put
$\Delta=\int\theta_T^6/6!$. Then
\[
 \boxed{\delta_W(A_T,\theta_T)
       \le10^{258}d^{69}f^{33}\Delta^6.}                  \tag{N9.1}
\]
Here $\delta_W$ is witnessed by a reduced irreducible effective
three-cycle with nonzero Weil projection. A sharper completely
specified bound is given by (N9.9) below.

\paragraph{Normalization with an explicit regularity integer.}
Keep $D=4d$, $r=64D$ and the sheaf $F$ of HD-N08. Set
\[
\begin{gathered}
 h=3\theta_D,\quad \mathcal D=\det F,\quad
 a=\frac{D-1}{8D},\\
 R=560r^2+8Dr+6=2294272D^2+6,\qquad
 B_0=[r]^*F\otimes\mathcal D^{-r}.
\end{gathered}                                             \tag{N9.2}
\]
The notation $B_0$ here denotes the ordinary normalized sheaf;
the rational B-field used below is written $\mathcal B_0$.
The class $h-\mathcal H$ is nef. Formula (N8.11) gives
$c_1(F)=-32DP+8(D-1)\theta_D$, and the class
$8Dh-c_1(F)$ has matrix
\[
 \begin{pmatrix}32D(2D+1)&32D\\32D&8(2D+1)\end{pmatrix},
 \qquad\det=256D(4D^2+1)>0.
\]
It is ample. For $1\le i\le6$, the line bundle
$\mathcal D^{-r}\otimes h^{R-i}$ therefore satisfies
\[
 c_1(\mathcal D^{-r}\otimes h^{R-i})-560r^2\mathcal H
 =560r^2(h-\mathcal H)+r(8Dh-c_1(F))+(6-i)h,
\]
which is ample. Assertion (N8.1) proves that $B_0$ is
$R$-regular. It still has projective dimension at most one.
Its complete character is
\[
 \operatorname{ch}(B_0)=e^{-r^2a\theta_D}
       \{-\lambda_Dr^6t_D+P_D(-r^2\theta_D)\}.            \tag{N9.3}
\]
In particular $c_1(B_0)=0$, and, writing its first two relevant
scalar coefficients relative to $h$,
\[
 \begin{aligned}
 \operatorname{ch}_2(B_0)&=a_2h^2,
 &a_2&=-\frac{r^4(D+1)^2}{18D},\\
 \operatorname{ch}_3(B_0)&=-\lambda_Dr^6t_D+a_3h^3,
 &a_3&=\frac{r^6(D-1)(D+1)^2}{648D^2}.
 \end{aligned}                                             \tag{N9.4}
\]

\paragraph{The precise deformation and ordinary-family comparison.}
Apply Perry's permitted theorem to the actual sheaf $F$, the
coherent finite action proved in HD-N08, and the algebraic
rational B-field $\mathcal B_0=-c_1(F)/r$. The whole twisted
character is the normalized expression in (N9.3) before the
$[r]$ scaling: it is a polynomial in the fixed polarization
plus a nonzero Weil class. It is flat and remains of Hodge type
through the balanced Weil family. The finite determinant
monodromy of the Weil line is killed as in HD-N07.

For precision about untwisting, Perry's Section6.2, proof of Theorem1.2,
printed pp.26--28, begins with the $\mu_n$-gerbe representing
the B-field denominator. Here take $n=r$. Its subsequent pushout
to $\mu_m$ does not change the associated Brauer class. On a
contractible analytic base the original class remains in the
local system $H^2(A,\mu_r)$. Multiplication by $r$ acts as
$r^2$ on this group and therefore kills it. Pullback by $[r]$
and an analytic Cech trivialization give an ordinary perfect
family. Here is its exact seed comparison, including the sign.
The seed gerbe is the $r$th-root gerbe of $\mathcal D^{-1}$;
its tautological weight-one line $\mathcal T_0$ satisfies
$\mathcal T_0^r=\mathcal D^{-1}$, and the seed object is
$F\otimes\mathcal T_0$. Write its class without an exponential
convention as the Kummer boundary
$\vartheta_0=\delta(\mathcal D^{-1})\in H^2(A,\mu_r)$.
Its rational B-field is
$c_1(\mathcal T_0)=c_1(\mathcal T_0^r)/r
=-c_1(\mathcal D)/r=\mathcal B_0$.
After pullback by $[r]$ and neutralization, the line becomes
an ordinary line $N_0$ with
$N_0^r=[r]^*\mathcal D^{-1}$. Hence
$c_1(N_0)=-r c_1(\mathcal D)$ and
$N_0=\mathcal D^{-r}\otimes P_0$ for some
$P_0\in\operatorname{Pic}^0(A)$. The ordinary seed is therefore
exactly $B_0\otimes P_0$, as asserted.
For the root gerbe, B-field and rational Chern character, see
Hotchkiss, \emph{Hodge theory of twisted derived categories and
the period-index problem}, Definition4.5 and Lemma4.6,
printed p.15; Remark4.13, printed p.17; Definition4.16 and
Theorem4.17, printed p.18,
in the original author version
\url{https://arxiv.org/pdf/2212.10638}.
The displayed root and tensor identities give the required
comparison directly. All preceding estimates allow the flat
factor $P_0$.

The locus where the derived fibres of a relative perfect
complex have cohomology only in degree zero is open. At the
seed this property holds by (N8.9); after shrinking the base
the ordinary family is therefore a coherent sheaf flat over
the base. This can equally be seen by the relative Tor-amplitude
criterion applied to a local finite vector-bundle resolution.
Properness permits a common neighbourhood for the whole fibre.
Analytic upper semicontinuity preserves the six regularity
vanishings of $B_0$ on a smaller neighbourhood. No stability
of all sheaves with this character is used.

\paragraph{A vector-bundle kernel supplies the effective cycle.}
At the seed and throughout that neighbourhood, evaluation gives
\[
 0\longrightarrow V\longrightarrow
    \mathcal O_A(-Rh)^N\longrightarrow B_0\longrightarrow0,
 \qquad N=h^0(B_0(Rh))=\chi(B_0(Rh)).                     \tag{N9.5}
\]
At the seed $V$ is locally free, because $\operatorname{pd}B_0\le1$.
This follows locally from the exact Ext sequence, or by comparing
the evaluation presentation with a projective resolution of
length one. In the family the direct image of $B_0(Rh)$ is a
rank-$N$ vector bundle, by the regularity vanishings and base
change; evaluation is surjective. Its kernel is flat, and its
locally free locus is open. The closed complement misses the
proper seed fibre, so it misses every fibre after one more
shrinking. Thus (N9.5) is an actual relative vector-bundle
presentation on the neighbourhood.

The kernel $V$ is $(R+7)$-regular. For cohomological degrees
at least two, use the sequence (N9.5), the regularity of $B_0$,
and the positive line bundles $(7-i)h$ for $2\le i\le6$.
For degree one, $R$-regularity gives the surjection
\[
 H^0(B_0(Rh))\otimes\operatorname{Sym}^6H^0(\mathbb P^M,\mathcal O(1))
 \longrightarrow H^0(B_0((R+6)h)),
\]
which factors through the degree-six sections in (N9.5).
It gives the remaining $H^1$ vanishing. Therefore
$V'=V((R+7)h)$ is a globally generated vector bundle.

Let $q_0=R+7$ and define the explicit rational numbers
\[
 \begin{split}
 u&=7N-rq_0,\\
 v&=\frac{49N}{2}-\frac{rq_0^2}{2}-a_2,\\
 w&=\frac{343N}{6}-\frac{rq_0^3}{6}-q_0a_2-a_3.
 \end{split}                                               \tag{N9.6}
\]
The actual identity in $K_0(A)$ is
\[
 [V']=N[\mathcal O_A(7h)]-[B_0(q_0h)].
\]
Consequently $c_1(V')=uh$, $\operatorname{ch}_2(V')=vh^2$,
and $\operatorname{ch}_3(V')=wh^3+\lambda_Dr^6t_D$.
The exact Newton identity
$c_3=2\operatorname{ch}_3-c_1\operatorname{ch}_2+c_1^3/6$
therefore gives
\[
 c_3(V')=2\lambda_Dr^6t_D+
                (2w-uv+u^3/6)h^3.                       \tag{N9.7}
\]
Its Weil projection is nonzero. In particular the rank of $V'$
is at least three, since lower-rank bundles have $c_3=0$.

General sections of a globally generated vector bundle give
an effective maximal-rank-drop locus of expected codimension
three, or the empty locus, representing $c_3(V')$. This is the
elementary Thom--Porteous construction: take
$\operatorname{rank}(V')-2$ sections, use the universal surjective
evaluation to get transversality to the rank strata, and then
take a general choice. Equation (N9.7) rules out the empty case.
It is an effective integral cycle with nonzero Weil projection,
so at least one reduced irreducible component detects the Weil
plane. No condition on the codimension of the nonfree locus of
$B_0$ is required.

Write $P_D(z)=\sum_{j=0}^6p_jz^j$, with coefficients given in
(N8.12). Since the Todd class of an abelian variety is1 and
$\theta_Dt_D=0$, the integer $N$ is explicitly
\[
 N=(6!\,64D^3)\sum_{j=0}^6
       p_j(-r^2)^j\frac{(3R-r^2a)^{6-j}}{(6-j)!}.
\]
The exact degree of the effective cycle above is the nonnegative
integer
\[
 C_D=(3^6\,6!\,64D^3)(2w-uv+u^3/6).                     \tag{N9.8}
\]
The formula defines every quantity by rational operations on
the specified positive integer $D$. Integrality and nonnegativity
also follow from the constructed vector bundle and effective
cycle. These same Chern classes and degrees persist on the
analytic neighbourhood. GAGA makes each fibre and cycle algebraic.

\paragraph{Actual source saturation and fixed-target isogenies.}
The source polarization $h=3\theta_D$ has product
$\Delta_h=3^6\,64D^3$. Since $D=4d$, the integral source order
$\mathbb Z[2\sqrt{-D}]=\mathbb Z[4\sqrt{-d}]$ has conductor
eight when $d\equiv3\pmod4$, and conductor four for $d=1,2$.
It contains the order of conductor eight in every case.
Let $\mathcal O$ be the maximal order and let
$\Lambda'=\mathcal O\Lambda$ be the saturation of the source
homology lattice. One has $8\Lambda'\subset\Lambda$.
For the Riemann form $E_h$, the form $8E_h$ is integral on
$\Lambda'$. To check this, write
$\mathcal O=\mathbb Z+\mathbb Z\alpha$: mixed terms use the
integral $8\alpha$ action, and
$E_h(\alpha x,\alpha y)=N(\alpha)E_h(x,y)$ follows from
Rosati conjugation. The actual quotient $q_s:A\to A'$ has
polarization $h'$ satisfying
\[
 q_s^*h'=8h,\qquad
 \Delta_s=\frac{8^6\Delta_h}{\deg q_s}
       \mid Q_*:=8^6\,3^6\,64D^3.
\]
Pushforward of the detecting cycle has degree at most $8^3C_D$
with respect to $h'$, and its Weil projection stays nonzero
under this $K$-linear isogeny.

For the nine fields in the statement, $\mathcal O$ is a PID.
Every split integral Hermitian lattice therefore has an actual
basis extending its primitive isotropic half-lattice. If its
polarization product is $q$, its matrix has form
$\left(\begin{smallmatrix}0&B\\B^*&D_0\end{smallmatrix}\right)$
with $N(\det B)=q$. Replacing
\[
 w\ \hbox{by}\ w-v(B^*)^{-1}D_0/2,\qquad
 u=w'B^{-1}
\]
gives the standard hyperbolic basis. Both this coordinate map
and its inverse are cleared by $2q$. In particular, the upper
right product $(B^*)^{-1}D_0B^{-1}$ is a block of the inverse
integral Hermitian matrix, whose determinant has absolute value
$q$; it does not require a denominator $q^2$.
The normalization is exact: for
$h_K(x,y)=E(x,\alpha y)-\bar\alpha E(x,y)$ the identity
Hermitian matrix induces the unimodular alternating matrix on
each integral basis $(v,\alpha v)$, so $|\det h_K|=q$.

Compare this basis for the source lattice with one for a fixed
maximal-order target lattice of product $\Delta$. The integer
$N_s=4Q_*\Delta$ clears the composition of the two rational
coordinate maps, since $\Delta_s\mid Q_*$. On the common
marked period domain this integral map induces an actual
isogeny $\phi:A'\to A_T$, with
$\phi^*\theta_T=N_s^2h'$. Passing to a common finite level
realizes it as a morphism of the corresponding abelian families.
It sends the nonempty analytic neighbourhood supplied by Perry
to an analytic open subset of the fixed target Weil component.
The pushforward degree cost is $N_s^6$. We have proved the
explicit bound
\[
 \boxed{\mathcal D_d(\Delta)=
       8^3(4Q_*\Delta)^6 C_{4d}.}                        \tag{N9.9}
\]
The bounded-detector closedness theorem of the supplied DEG\_L
framework then extends the same bound to every fibre of that
fixed polarized component. This uses a fixed degree bound
already proved on a nonempty analytic open, rather than an
unbounded collection of isogenies.

\paragraph{A simpler explicit bound.}
Put $s=3R-r^2a$. Directly from (N8.12),
$j!|p_j|\le64D$ for $0\le j\le6$. Also
$s+r^2<7\cdot10^6D^2$. The exact Euler-characteristic formula
therefore gives
\[
 0<N\le4096D^4(s+r^2)^6
       <10^{45}D^{16}.
\]
Here $q_0<3\cdot10^6D^2$,
$|a_2|<4\cdot10^6D^5$ and $|a_3|<5\cdot10^8D^7$.
The formulas (N9.6) imply
\[
 |u|\le8\cdot10^{45}D^{16},\quad
 |v|\le26\cdot10^{45}D^{16},\quad
 |w|\le59\cdot10^{45}D^{16}.
\]
Since $10^{45}D^{16}\ge1$, the expression in (N9.8) is bounded
in absolute value by $512\cdot10^{135}D^{48}$. Hence
\[
 0<C_D<10^{146}D^{51},\qquad C_{4d}<10^{177}d^{51}.
\]
Moreover
$Q_*=782757789696\,d^3<10^{12}d^3$, and therefore
$8^3(4Q_*\Delta)^6<10^{81}d^{18}\Delta^6$.
Together these inequalities give
$\mathcal D_d(\Delta)<10^{258}d^{69}\Delta^6$.

Finally, for a target order of conductor $f$, saturate the
target lattice to $\mathcal O\Lambda$ and use the integral
Riemann form $fE$. If the quotient has degree $b$, then
$q_f^*\theta'=f\theta_T$,
$\Delta'=f^6\Delta/b$, and pullback of a three-cycle costs
$b/f^3$. The degree-six homogeneity of (N9.9) gives the exact
factor $f^{33}/b^5$, bounded by $f^{33}$. The same factor applies
to the coarse bound, proving (N9.1). \hfill$\square$

\paragraph{Provenance and check scope.}
The original secant-square design is credited to Markman's supplied
preprint. HD-N08 supplies the actual quotient ideals, semiregularity,
Fourier morphism, sheaf property and numerical regularity directly;
it does not take that preprint's semiregularity or local-freeness
theorems as inputs. The regularity method follows the explicit
three-term and Buchsbaum--Rim calculation in the attached
HECKE\_DIV M7, with every rank, twist and positivity matrix changed
and checked here. The new evaluation-kernel construction replaces
its codimension-four nonfree-locus requirement by the proved
projective-dimension-one presentation. The deformation input is
the permitted Perry theorem, with the stated actual gerbe comparison.
